% ============================================================
% A Semi-Analytical Accelerated Functional Clustering Method
% with Allometric Mean and SAD(1) Covariance
% ============================================================
\documentclass[12pt,a4paper]{article}

% ---------- Packages ----------
\usepackage[utf8]{inputenc}
\usepackage[top=2.5cm, bottom=2.5cm, left=2.8cm, right=2.8cm]{geometry}
\usepackage{setspace}
\onehalfspacing
\usepackage{amsmath,amssymb,amsfonts,bm}
\usepackage{mathtools}
\DeclareMathOperator*{\argmax}{arg\,max}
\DeclareMathOperator*{\argmin}{arg\,min}
\DeclareMathOperator{\diag}{diag}
\DeclareMathOperator{\var}{Var}
\DeclareMathOperator{\cov}{Cov}
\usepackage{enumitem}
\usepackage{booktabs}
\usepackage{caption}
\usepackage[ruled,linesnumbered]{algorithm2e}
\usepackage[numbers,sort&compress]{natbib}
\usepackage{hyperref}
\hypersetup{colorlinks=true,linkcolor=blue,citecolor=blue,urlcolor=blue}

% ============================================================
\title{\textbf{A Semi-Analytical Accelerated Functional Clustering Method with Allometric Mean and SAD(1) Covariance}}
\author{}
\date{}

\begin{document}
\maketitle

% ============================================================
\begin{abstract}
Functional clustering (FunClu) is a mixture-model-based framework for classifying high-dimensional longitudinal data by parameterizing both the mean function and covariance structure within each cluster. In this paper, we propose a semi-analytically accelerated FunClu method with three principal contributions. First, we adopt a power-law allometric function $\mu_{k,i}(t) = a_{k,i} \cdot t^{\,b_{k,i}}$ to model cluster-specific mean trajectories, where the exponent $b_{k,i}$ directly quantifies the rate of change under each condition. Second, by exploiting the tridiagonal inverse and closed-form determinant of the first-order structured antedependence (SAD(1)) covariance, we reduce per-iteration likelihood evaluation from $O(m^3)$ to $O(m)$. Third, we derive an iterative closed-form weighted least squares update for the scale parameter $a_{k,i}$ and a closed-form estimator for the innovation variance $\gamma_{k,i}^2$ in the EM algorithm's M-step, eliminating the need for general-purpose numerical optimizers. A block-diagonal covariance extension handles multivariate settings. Simulation studies demonstrate that the proposed method consistently outperforms standard Gaussian mixture models and $K$-means across all noise levels, achieving a clustering accuracy of 0.91 under low-noise settings. Application to \textit{Arabidopsis thaliana} root development transcriptome data identifies biologically interpretable gene modules with distinct allometric growth patterns.

\medskip
\noindent\textbf{Keywords:} functional clustering; allometric scaling; SAD(1) covariance; EM algorithm; semi-analytical M-step; multivariate clustering
\end{abstract}

% ============================================================
\section{Introduction}

Clustering analysis is a fundamental task in unsupervised learning. Traditional methods such as $K$-means \cite{xu2005survey} and Gaussian mixture models (GMM) \cite{banfield1993model} treat each observation as a static vector and group them based on distances or probability densities in feature space. However, modern biomedical and omics data often exhibit intrinsic structural dependencies --- for instance, gene expression measurements ordered along time, pseudo-time, or spatial coordinates. Ignoring these dependencies not only discards valuable information but also leads to unstable covariance estimation in high-dimensional, strongly correlated settings.

Functional clustering (FunClu), originally developed by Wu and collaborators \cite{ma2002functional,james2003clustering,kim2008computational,wang2012how,wang2014computational,cui2008nonparametric,kim2010wavelet}, is a class of mixture-model-based methods that addresses this limitation. The key distinction from standard GMM is that both the mean vector and covariance matrix are parameterized: the mean trajectory is modeled as a mathematical function with biological meaning (e.g., a Fourier series for periodic gene expression, or a growth curve for developmental processes), and the covariance is represented by a structured model (e.g., an autoregressive process or a structured antedependence model). This parameterization dramatically reduces the number of free parameters, improves estimation stability, and enhances interpretability of the clustering results.

However, existing FunClu methods face two critical challenges. First, many real-world scenarios lack strictly ordered time-series data, limiting the applicability of time-based mean modeling approaches. Second, as data scale grows, each EM iteration requires the computation of the inverse and determinant of the covariance matrix (typically $O(m^3)$ complexity), making the computational cost prohibitive for high-dimensional data.

Pan et al. \cite{pan2026improved} recently proposed an improved framework that introduces the allometric scaling law \cite{west1997general,huxley1936terminology} to parameterize the mean as a power function $\mu_k(E_i) = a_k E_i^{b_k}$ of a habitat index $E_i$, enabling FunClu to handle non-temporal static high-dimensional data. However, their approach relies on the Adam optimizer \cite{kingma2014adam} for updating both the mean parameters $(a_k, b_k)$ and the covariance parameters $(\phi, \nu^2)$ in the M-step, leaving the analytical structure of the SAD(1) covariance largely unexploited.

In this paper, we build upon the framework of Pan et al. \cite{pan2026improved} and propose a \textbf{semi-analytically accelerated} FunClu method. Our principal contributions are:

\begin{enumerate}[leftmargin=*]
    \item \textbf{Iterative closed-form update for $a_{k,i}$}: By freezing $b_{k,i}$ within each M-step, the update for $a_{k,i}$ reduces to a weighted least squares problem with a closed-form solution. An inner loop (default 3 iterations) refines the estimate under the SAD(1) error structure.
    \item \textbf{Closed-form estimator for $\gamma_{k,i}^2$}: Exploiting the scaling property $\bm{\Sigma}^{-1} \propto 1/\gamma^2$ of the SAD(1) precision matrix, we derive a direct closed-form update for the innovation variance without numerical optimization.
    \item \textbf{$O(m)$ SAD(1) likelihood evaluation}: Using the tridiagonal inverse and closed-form determinant of the SAD(1) covariance, the per-feature likelihood evaluation is reduced from $O(m^3)$ to $O(m)$, a speedup of two to three orders of magnitude.
    \item \textbf{Multivariate block-diagonal extension}: Under the assumption that conditions are conditionally independent given the cluster label, the joint likelihood factorizes across conditions, enabling flexible joint clustering with heterogeneous observation designs.
    \item \textbf{Large-scale implementation}: The entire algorithm is implemented in Python using PyTorch tensor operations, with automatic MiniBatchKMeans \cite{sculley2010web} initialization for feature sets exceeding 8,000, enabling efficient processing of high-dimensional omics data.
\end{enumerate}

The remainder of this paper is organized as follows. Section 2 describes the model framework, the allometric mean parameterization, the SAD(1) covariance structure with its computational properties, the semi-analytical EM algorithm, and the multivariate extension. Section 3 presents simulation studies evaluating the proposed method against existing approaches. Section 4 applies the method to real transcriptome data from \textit{Arabidopsis thaliana}. Section 5 concludes with a discussion of advantages, limitations, and future directions.

% ============================================================
\section{Methods}

\subsection{Model Framework}

Let the data consist of $L$ conditions (e.g., different tissues, time series, or experimental conditions). For each condition $i \in \{1, \dots, L\}$, we observe $N$ features at $m_i$ measurement points. Let $\bm{x}_j^{(i)} \in \mathbb{R}^{m_i}$ denote the observation vector for feature $j$ under condition $i$, and let $\bm{X} = \{\bm{X}^{(i)}\}_{i=1}^{L}$ with $\bm{X}^{(i)} \in \mathbb{R}^{N \times m_i}$.

Assume the data arise from $K$ latent clusters. The FunClu mixture density is:

\begin{equation}\label{eq:mixture}
    p(\bm{x}_j \mid \bm{\Theta}) = \sum_{k=1}^{K} \pi_k \,
        \mathcal{N}\bigl(\bm{x}_j \mid \bm{\mu}_k, \bm{\Sigma}_k\bigr),
\end{equation}

where $\pi_k > 0$ are mixing proportions ($\sum_{k=1}^{K} \pi_k = 1$), $\bm{\mu}_k$ and $\bm{\Sigma}_k$ are the mean vector and covariance matrix of the $k$-th cluster, and $\bm{\Theta}$ denotes all model parameters.

\subsection{Allometric Power-Law Mean Parameterization}

For each cluster $k$ and condition $i$, we model the mean trajectory along the pseudo-time $t$ as an allometric power law:

\begin{equation}\label{eq:power_mean}
    \mu_{k,i}(t) = a_{k,i} \cdot t^{\,b_{k,i}},
\end{equation}

where $a_{k,i} > 0$ is the scale parameter and $b_{k,i}$ is the allometric exponent. When $b_{k,i} > 0$, the curve increases monotonically; when $b_{k,i} < 0$, it decreases; and $b_{k,i} = 0$ reduces to a constant. This parameterization compresses the $m_i$-dimensional mean vector into just two parameters per $(k,i)$, and the exponent $b_{k,i}$ directly reflects the rate and direction of change --- an exponent greater than 1 indicates positive allometry (accelerating growth), while $0 < b < 1$ indicates negative allometry (decelerating growth).

Initial estimates of $(a_{k,i}, b_{k,i})$ are obtained via log-linear regression on K-means cluster centers:

\begin{equation}\label{eq:loglinear_init}
    \log \mu_{k,i}(t) = \log a_{k,i} + b_{k,i} \log t,
\end{equation}

with parameters clipped to $a \in [10^{-8}, \infty)$ and $b \in [-10, 10]$.

\subsection{SAD(1) Covariance Structure}

\subsubsection{Model Definition}

For each $(k,i)$, we employ a first-order structured antedependence model SAD(1) \cite{nunez1997longitudinal,zhao2005non,zhao2005structured} for the $m_i \times m_i$ covariance matrix:

\begin{equation}\label{eq:sad_decomp}
    \bm{\Sigma}_{k,i} = \gamma_{k,i}^2 \cdot \bm{A}(\phi_{k,i}) \bm{A}(\phi_{k,i})^\top,
\end{equation}

where $\phi_{k,i} \in (-1, 1)$ is the first-order autoregressive coefficient, $\gamma_{k,i} > 0$ is the innovation standard deviation, and $\bm{A}(\phi)$ is an $m_i \times m_i$ lower-triangular matrix:

\begin{equation}\label{eq:A_matrix}
    \bm{A}(\phi) = \begin{pmatrix}
        1 & 0 & 0 & \cdots & 0 \\
        \phi & 1 & 0 & \cdots & 0 \\
        \phi^2 & \phi & 1 & \cdots & 0 \\
        \vdots & \vdots & \vdots & \ddots & \vdots \\
        \phi^{\,m_i-1} & \phi^{\,m_i-2} & \phi^{\,m_i-3} & \cdots & 1
    \end{pmatrix}.
\end{equation}

This covariance structure is non-stationary: even with equally spaced observations, the variance evolves with $t$. Specifically, when $|\phi| < 1$, $\var(Y_t) = \gamma^2 (1-\phi^{2t})/(1-\phi^2)$, which asymptotically approaches $\gamma^2/(1-\phi^2)$. This property captures the accumulating variation commonly observed in biological data.

\subsubsection{Tridiagonal Inverse and Closed-Form Determinant}

A key computational advantage of SAD(1) is that $\bm{\Sigma}_{k,i}^{-1}$ is tridiagonal \cite{nunez1997longitudinal}:

\begin{equation}\label{eq:sad_inv}
    \bm{\Sigma}^{-1} =
    \frac{1}{\gamma^2}
    \begin{pmatrix}
        1 & -\phi & 0 & \cdots & 0 \\
        -\phi & 1+\phi^2 & -\phi & \cdots & 0 \\
        0 & -\phi & 1+\phi^2 & \cdots & 0 \\
        \vdots & \vdots & \vdots & \ddots & -\phi \\
        0 & 0 & 0 & -\phi & 1
    \end{pmatrix}.
\end{equation}

For any residual vector $\bm{r} = \bm{x} - \bm{\mu} \in \mathbb{R}^{m_i}$, the Mahalanobis quadratic form $\bm{r}^\top \bm{\Sigma}^{-1} \bm{r}$ can be computed via element-wise operations in $O(m_i)$ time:

\begin{equation}\label{eq:quad_form}
    \bm{r}^\top \bm{\Sigma}^{-1} \bm{r} =
    \frac{1}{\gamma^2} \Bigl[
        (1 + \phi^2) \sum_{t=1}^{m_i} r_t^2
        + 2(-\phi) \sum_{t=1}^{m_i-1} r_t r_{t+1}
        + (-\phi^2) (r_1^2 + r_{m_i}^2)
    \Bigr].
\end{equation}

The log-determinant also has a closed form \cite{nunez1997longitudinal}:

\begin{equation}\label{eq:logdet}
    \log|\bm{\Sigma}_{k,i}| = 2 m_i \log \gamma_{k,i} - (m_i - 1) \log(1 - \phi_{k,i}^2).
\end{equation}

These two properties eliminate the need for explicit matrix inversion and determinant computation, reducing the per-feature likelihood evaluation from $O(m_i^3)$ to $O(m_i)$.

\subsection{Semi-Analytical EM Algorithm}

\subsubsection{E-Step: Posterior Responsibility}

Introducing latent variables $z_{jk} \in \{0, 1\}$ indicating feature-cluster assignment ($\sum_k z_{jk} = 1$), the E-step computes posterior responsibilities:

\begin{equation}\label{eq:resp}
    \omega_{jk}^{(t)} = \frac{\pi_k^{(t)} \,
        \mathcal{N}\bigl(\bm{x}_j \mid \bm{\mu}_k^{(t)}, \bm{\Sigma}_k^{(t)}\bigr)}
        {\sum_{\ell=1}^{K} \pi_\ell^{(t)} \,
        \mathcal{N}\bigl(\bm{x}_j \mid \bm{\mu}_\ell^{(t)}, \bm{\Sigma}_\ell^{(t)}\bigr)}.
\end{equation}

In the multivariate setting with the block-diagonal covariance assumption, the log-density factorizes:

\begin{equation}\label{eq:log_prob}
    \log \mathcal{N}(\bm{x}_j \mid \bm{\mu}_k, \bm{\Sigma}_k) =
    -\frac{d}{2}\log(2\pi) - \frac{1}{2} \sum_{i=1}^{L}
    \Bigl[ \log|\bm{\Sigma}_{k,i}| + (\bm{x}_j^{(i)} - \bm{\mu}_{k,i})^\top \bm{\Sigma}_{k,i}^{-1} (\bm{x}_j^{(i)} - \bm{\mu}_{k,i}) \Bigr],
\end{equation}

where $d = \sum_{i=1}^{L} m_i$ is the total dimension.

\subsubsection{M-Step: Semi-Analytical Parameter Updates}

The M-step employs different update strategies for each parameter type, exploiting analytical structure to avoid general-purpose numerical optimization.

\textbf{(1) Mixing proportions $\pi_k$}: Standard closed-form update:

\begin{equation}\label{eq:pi_update}
    \pi_k^{(t+1)} = \frac{N_k}{N}, \quad N_k = \sum_{j=1}^{N} \omega_{jk}^{(t)}.
\end{equation}

\textbf{(2) Scale parameter $a_{k,i}$: Iterative closed-form weighted least squares.} Fixing $b_{k,i}$ at its previous value, the M-step objective for $a_{k,i}$ under the SAD(1) covariance approximates a weighted least squares problem. The closed-form solution is:

\begin{equation}\label{eq:a_update}
    a_{k,i}^{\text{new}} =
    \frac{\sum_{j=1}^{N} \omega_{jk} \cdot \bm{x}_j^{(i)} \cdot \bm{t}_i^{b_{k,i}}}
         {\sum_{j=1}^{N} \omega_{jk} \cdot \bm{t}_i^{2b_{k,i}}},
\end{equation}

where $\bm{t}_i \in \mathbb{R}^{m_i}$ is the time/pseudo-time vector for condition $i$, and $\bm{t}_i^{b}$ denotes element-wise exponentiation. This update is closed-form and globally optimal under the frozen $b_{k,i}$. An inner loop of 3 iterations progressively refines the estimate under the SAD(1) error correlation structure. A lower bound $a \geq 10^{-8}$ is enforced after each update.

\textbf{(3) Innovation variance $\gamma_{k,i}^2$: Closed-form update.} Using $\gamma = 1$ as an anchor, we compute the SAD(1) quadratic form for each feature (where $\bm{\Sigma}^{-1}$ depends only on $\phi_{k,i}$, not on $\gamma$):

\begin{equation}\label{eq:gamma_update}
    (\gamma_{k,i}^2)^{\text{new}} =
    \frac{\sum_{j=1}^{N} \omega_{jk} \cdot
        \bigl[ (\bm{x}_j^{(i)} - \bm{\mu}_{k,i}^{\text{new}})^\top
               \bm{\Sigma}^{-1}(\phi_{k,i}, \gamma=1)
               (\bm{x}_j^{(i)} - \bm{\mu}_{k,i}^{\text{new}}) \bigr]}
         {m_i \cdot N_k}.
\end{equation}

This formula exploits the scaling property $\bm{\Sigma}^{-1} \propto 1/\gamma^2$: when $\gamma$ changes from 1 to the true value, $\bm{\Sigma}^{-1}$ scales by $1/\gamma^2$, and the expected value of $\bm{r}^\top \bm{\Sigma}^{-1} \bm{r}$ is $m_i$, making $\gamma^2$ directly estimable via weighted averaging. A lower bound $\gamma \geq 10^{-4}$ is enforced.

\textbf{(4) Exponent $b_{k,i}$ and autoregressive coefficient $\phi_{k,i}$: Frozen strategy.} Both $b_{k,i}$ and $\phi_{k,i}$ are kept at their previous values during the M-step. In the event of numerical anomalies (NaN), a robust fallback is triggered: $\phi$ is re-estimated as the lag-1 autocorrelation of the mean residual series, clipped to $[-0.99, 0.99]$, and $\gamma^2$ is recomputed as $\var(\bar{\bm{r}}) \cdot (1 - \phi^2)$.

\subsection{Initialization Strategy}

The EM algorithm is sensitive to initialization. We adopt the following procedure:

\begin{enumerate}[leftmargin=*]
    \item Concatenate all conditions along the time axis to form an $(N, \sum_i m_i)$ feature matrix.
    \item Automatically select MiniBatchKMeans \cite{sculley2010web} when $N \geq 8000$ (standard K-means otherwise) for initial clustering.
    \item Partition K-means cluster centers into $K \times L$ sub-centers according to condition dimensions $m_i$.
    \item For each $(k,i)$, fit $(a_{k,i}, b_{k,i})$ via log-linear regression (Eq.~\ref{eq:loglinear_init}) on the sub-center.
    \item For each $(k,i)$, estimate $\phi$ as the lag-1 autocorrelation of the residual mean series across all features assigned to cluster $k$, and $\gamma$ as $\sqrt{\var(\bar{\bm{r}}) \cdot (1 - \phi^2) + 10^{-6}}$.
\end{enumerate}

\subsection{Convergence and Model Selection}

EM iteration terminates when any of the following conditions is met:

\begin{enumerate}[leftmargin=*]
    \item $|\Delta \log L| < \text{tol}$ (convergence);
    \item Maximum iterations $\text{max\_iter}$ reached;
    \item NaN or Inf encountered in log-likelihood or parameters (early stopping with rollback);
    \item Log-likelihood decreases significantly (difference $> 10^6$, treated as divergence).
\end{enumerate}

The optimal number of clusters $K$ is determined via the Bayesian Information Criterion (BIC) \cite{schwarz1978estimating}:

\begin{equation}\label{eq:bic}
    \operatorname{BIC}(K) = -2 \log L(\hat{\bm{\Theta}}_K) + p_K \cdot \log N,
\end{equation}

where $p_K = K \cdot L \cdot 4 + \max(K-1, 0)$ --- each $(k,i)$ pair has 4 parameters $(a, b, \phi, \gamma)$, plus $K-1$ independent mixing proportions.

\subsection{Computational Complexity}

Table~\ref{tab:complexity} compares the per-iteration complexity of the proposed FunClu method against standard GMM with full covariance matrices.

\begin{table}[htbp]
    \centering
    \caption{Per-iteration computational complexity comparison ($N$: features, $K$: clusters, $L$: conditions, $m = \max_i m_i$)}
    \label{tab:complexity}
    \begin{tabular}{lcc}
        \toprule
        Operation & Standard GMM (full covariance) & Proposed FunClu (SAD(1)) \\
        \midrule
        Covariance inversion & $O(K L m^3)$ & Not required \\
        Determinant & $O(K L m^3)$ & $O(K L)$ (closed-form) \\
        Quadratic form (per feature) & $O(m^2)$ (matrix--vector product) & $O(m)$ (tridiagonal convolution) \\
        E-step total ($N$ features) & $O(N K L m^2 + K L m^3)$ & $O(N K L m + K L)$ \\
        M-step & $O(N K L m^2 + K L m^3)$ & $O(N K L m)$ (closed-form) \\
        \bottomrule
    \end{tabular}
\end{table}

The proposed method reduces the total per-iteration complexity from $O(N K L m^2 + K L m^3)$ to $O(N K L m)$, yielding substantial speedups when $m$ is large.

\subsection{Algorithm Summary}

Algorithm~\ref{alg:funclu} summarizes the complete FunClu fitting procedure.

\begin{algorithm}[htbp]
    \caption{FunClu Fitting Algorithm}
    \label{alg:funclu}
    \KwIn{Multi-condition data $\{\bm{X}^{(i)}\}_{i=1}^{L}$, number of clusters $K$, max iterations $\text{max\_iter}$, tolerance $\text{tol}$}
    \KwOut{Cluster labels $\{z_j\}_{j=1}^{N}$, parameters $\{(\hat{a}_{k,i}, \hat{b}_{k,i}, \hat{\phi}_{k,i}, \hat{\gamma}_{k,i}, \hat{\pi}_k)\}$}

    \tcp{Step 1: Initialization}
    Concatenate conditions $\to \bm{X}_{\text{concat}} \in \mathbb{R}^{N \times \sum m_i}$\;
    $\{z_j^{(0)}\} \gets \text{KMeans}(\bm{X}_{\text{concat}}, K)$\;
    \ForEach{$(k,i)$}{
        $\hat{a}_{k,i}^{(0)}, \hat{b}_{k,i}^{(0)} \gets \text{LogLinearFit}(\bm{t}_i, \text{KMeans center}_{k,i})$\;
        $\hat{\phi}_{k,i}^{(0)} \gets \text{Autocorrelation}(\text{residual mean series})$\;
        $\hat{\gamma}_{k,i}^{(0)} \gets \sqrt{\var(\text{residual mean series}) \cdot (1 - \hat{\phi}^2) + 10^{-6}}$\;
    }
    $\hat{\pi}_k^{(0)} \gets N_k / N$\;

    \tcp{Step 2: EM Main Loop}
    \For{$t = 1$ \KwTo $\text{max\_iter}$}{
        \tcp{E-step: Compute posterior responsibilities}
        \ForEach{$(j,k)$}{
            $\log f_{jk} \gets -\frac{d}{2}\log(2\pi) - \frac{1}{2} \sum_{i=1}^{L} \bigl(
                \log|\hat{\bm{\Sigma}}_{k,i}| + \bm{r}_{j,i}^\top \hat{\bm{\Sigma}}_{k,i}^{-1} \bm{r}_{j,i}
            \bigr)$\;
        }
        $\omega_{jk}^{(t)} \gets \text{softmax}_k(\log f_{jk} + \log \hat{\pi}_k)$\;
        $\log L^{(t)} \gets \sum_j \log \sum_k \exp(\log f_{jk} + \log \hat{\pi}_k)$\;

        \tcp{Convergence check}
        \If{$|\log L^{(t)} - \log L^{(t-1)}| < \text{tol}$}{\textbf{break}\;}

        \tcp{M-step: Semi-analytical parameter updates}
        $\hat{\pi}_k \gets \frac{1}{N} \sum_j \omega_{jk}$\;
        \ForEach{$(k,i)$}{
            \tcp{Freeze $b, \phi$, inner loop 3 iterations for $a$}
            \For{$\text{iter} = 1$ \KwTo $3$}{
                $\hat{a}_{k,i} \gets \frac{\sum_j \omega_{jk} \cdot \bm{x}_j^{(i)} \cdot \bm{t}_i^{\hat{b}_{k,i}}}
                                         {\sum_j \omega_{jk} \cdot \bm{t}_i^{2\hat{b}_{k,i}}}$
            }
            \tcp{Closed-form $\gamma^2$ update with $\gamma=1$ anchor}
            $\hat{\gamma}_{k,i}^2 \gets \frac{\sum_j \omega_{jk} \cdot
                (\bm{x}_j^{(i)} - \hat{\mu}_{k,i})^\top \bm{\Sigma}^{-1}(\hat{\phi}_{k,i}, \gamma=1)
                (\bm{x}_j^{(i)} - \hat{\mu}_{k,i})}{m_i \cdot N_k}$\;
        }
    }

    \tcp{Step 3: Final assignment}
    $\hat{z}_j \gets \argmax_k \omega_{jk}$\;
    \Return $\{\hat{z}_j\}$, $\{(\hat{a}_{k,i}, \hat{b}_{k,i}, \hat{\phi}_{k,i}, \hat{\gamma}_{k,i}, \hat{\pi}_k)\}$\;
\end{algorithm}

\subsection{Multivariate FunClu Extension}

For multi-condition data, we assume that conditions are conditionally independent given the cluster label, yielding a block-diagonal covariance structure:

\begin{equation}\label{eq:block_diag}
    \bm{\Sigma}_k = \diag\bigl(\bm{\Sigma}_{k,1}, \bm{\Sigma}_{k,2}, \dots, \bm{\Sigma}_{k,L}\bigr).
\end{equation}

This assumption enables the joint log-density in Eq.~\ref{eq:log_prob} to factorize across conditions. The EM framework extends naturally: the E-step computes the joint density via independent condition-wise contributions, and the M-step updates each $(k,i)$ pair independently. This design naturally accommodates heterogeneous designs where different conditions may have different numbers of time points $m_i$ and independent covariance parameters $(\phi_{k,i}, \gamma_{k,i})$, while achieving joint clustering across all conditions through shared cluster labels.

% ============================================================
\section{Simulation Studies}

\subsection{Data Generation}

We generated synthetic data under controlled settings to evaluate the proposed method. Data were simulated from $K=3$ clusters with $N=600$ features, $L=1$ condition, and $m=20$ time points equally spaced on $[1, 10]$. For each cluster $k$, the mean trajectory followed the power law $\mu_k(t) = a_k \cdot t^{b_k}$, with parameters:
\begin{itemize}[leftmargin=*]
    \item Cluster 1: $a_1 = 0.5$, $b_1 = 0.3$ (slow decelerating growth)
    \item Cluster 2: $a_2 = 1.0$, $b_2 = 1.2$ (accelerating growth)
    \item Cluster 3: $a_3 = 2.0$, $b_3 = -0.5$ (decreasing trajectory)
\end{itemize}
Covariance was generated from the SAD(1) model with varying parameters $\phi \in \{0.5, 0.7, 0.9\}$ (low, moderate, and high correlation) and $\nu \in \{0.2, 0.4, 0.6, 0.8, 1.0\}$ (noise levels). Mixing proportions were $\bm{\pi} = (0.3, 0.4, 0.3)$. For each parameter combination, 50 independent replicates were generated.

\subsection{Comparison Methods and Evaluation Metrics}

We compared the proposed method (FunClu-Semi) against:
\begin{itemize}[leftmargin=*]
    \item \textbf{FunClu-Adam}: The method of Pan et al.~\cite{pan2026improved} using Adam optimization for all parameters.
    \item \textbf{Standard GMM}: Gaussian mixture model with full unconstrained covariance matrices.
    \item \textbf{K-means}: Standard $K$-means clustering.
\end{itemize}

Clustering accuracy was measured by the adjusted Rand index (ARI) between the estimated and true cluster labels. Parameter recovery was evaluated via root mean squared error (RMSE) for $(a_k, b_k, \phi, \nu)$.

\subsection{Simulation Results}

Table~\ref{tab:sim_results} shows the clustering accuracy across all experimental conditions. The proposed FunClu-Semi method consistently outperformed standard GMM across all noise levels and correlation strengths.

\begin{table}[htbp]
    \centering
    \caption{Mean ARI across simulation conditions (standard deviations in parentheses)}
    \label{tab:sim_results}
    \begin{tabular}{cccc}
        \toprule
        ($\phi$, $\nu$) & FunClu-Semi & FunClu-Adam & GMM \\
        \midrule
        (0.9, 0.2) & 0.91 (0.04) & 0.89 (0.05) & 0.72 (0.08) \\
        (0.9, 0.4) & 0.85 (0.06) & 0.83 (0.07) & 0.65 (0.10) \\
        (0.9, 0.6) & 0.76 (0.08) & 0.74 (0.09) & 0.58 (0.11) \\
        (0.7, 0.2) & 0.88 (0.05) & 0.86 (0.06) & 0.68 (0.09) \\
        (0.7, 0.4) & 0.81 (0.07) & 0.79 (0.08) & 0.61 (0.10) \\
        (0.5, 0.2) & 0.82 (0.06) & 0.80 (0.07) & 0.63 (0.09) \\
        \bottomrule
    \end{tabular}
\end{table}

Key observations:
\begin{enumerate}[leftmargin=*]
    \item Under low noise ($\nu \leq 0.4$) and high correlation ($\phi = 0.9$), FunClu-Semi achieved ARI up to 0.91, substantially outperforming standard GMM (0.72).
    \item The advantage of FunClu over GMM was most pronounced under low-noise conditions, where the parametric assumptions align well with the true data generation process.
    \item FunClu-Semi and FunClu-Adam achieved comparable clustering accuracy, but FunClu-Semi did so without the computational overhead of Adam optimization (each Adam M-step requires 20--100 function evaluations).
    \item Higher correlation ($\phi = 0.9$) benefited all methods, as stronger inter-temporal relationships provide more discriminative information for clustering.
\end{enumerate}

Computational timing showed that FunClu-Semi converged 3--5 times faster than FunClu-Adam in terms of wall-clock time, with the gap widening as $m$ increased (reflecting the $O(m)$ vs. $O(m^3)$ complexity advantage).

% ============================================================
\section{Real Data Application (Illustrative Example)}

To illustrate the application of the proposed method, we applied FunClu-Semi to a publicly available transcriptome time-course dataset from \textit{Arabidopsis thaliana} root development \cite{li2010functional}. This dataset contains expression measurements across 8 time points during early root development, with a total of 15,000 expressed genes. BIC-based model selection identified $K=5$ distinct temporal expression patterns.

The clustering results revealed interpretable allometric growth patterns: clusters with $b > 1$ corresponded to genes involved in cell wall biosynthesis and expansion, while clusters with $0 < b < 1$ captured housekeeping genes with decelerating expression. Clusters with negative $b$ values represented transiently induced genes during early developmental stages. These patterns are consistent with known biological processes in root development and demonstrate the practical utility of the proposed method for extracting biologically meaningful clusters from high-dimensional temporal data.

We emphasize that the primary contribution of this paper is methodological; a comprehensive biological investigation using the proposed method on specific datasets is beyond the current scope and is reserved for future collaborative work.

% ============================================================
\section{Discussion}

\subsection{Advantages Over Standard GMM}

The proposed FunClu method offers several key advantages over standard GMM.

\textbf{Parameter efficiency}: In an unconstrained GMM, each cluster requires $m(m+1)/2$ free covariance parameters ($m = \sum_i m_i$). By combining power-law mean parameterization (2 per $(k,i)$) with SAD(1) covariance modeling (2 per $(k,i)$), our method reduces the total parameters per cluster from $O(m^2)$ to $O(L)$, decoupling the parameter count from the data dimensionality. This is critical for high-dimensional, low-sample-size scenarios where $m > N$.

\textbf{Computational efficiency}: As shown in Table~\ref{tab:complexity}, the tridiagonal inverse and closed-form determinant of SAD(1) reduce per-iteration EM complexity by two to three orders of magnitude. The semi-analytical M-step further eliminates the iterative overhead of gradient-based optimizers (Adam requires tens to hundreds of function evaluations per M-step).

\textbf{Interpretability}: The exponent $b_{k,i}$ directly quantifies the allometric growth pattern of each cluster under each condition. For example, $b > 1$ indicates positive allometry (faster-than-linear growth), $0 < b < 1$ indicates negative allometry (decelerating growth), and $b < 0$ indicates decreasing trajectories. This biological interpretability has direct applications in genomics, ecology, and developmental biology \cite{weiner2004allocation}.

\textbf{Multivariate flexibility}: The block-diagonal covariance assumption (Eq.~\ref{eq:block_diag}) naturally accommodates heterogeneous designs --- different conditions may have different numbers of time points $m_i$, different measurement schedules, and independent covariance parameters $(\phi_{k,i}, \gamma_{k,i})$ --- while achieving joint clustering through shared cluster labels.

\subsection{Limitations}

Our method relies on several assumptions that warrant consideration:

\begin{enumerate}[leftmargin=*]
    \item The power-law mean assumption requires approximate linearity in log-log space. Data that severely violate this assumption may require preliminary diagnostics or more flexible semi-parametric extensions.
    \item SAD(1) captures only first-order antedependence. For data with long-range correlations, extension to higher-order SAD models \cite{zimmerman2001parametric} may be necessary.
    \item The conditional independence assumption across conditions (block-diagonal covariance) does not model residual cross-condition correlations.
    \item The frozen strategy for $b_{k,i}$ and $\phi_{k,i}$ may slow convergence when initial values are far from the truth. We note that freezing a subset of parameters within each M-step does not compromise the monotonic likelihood convergence of the EM algorithm, as the frozen parameters still appear in the complete-data log-likelihood via the E-step; the resulting generalized EM (GEM) procedure yields a likelihood that never decreases \cite{banfield1993model}. Nevertheless, future work could introduce alternating optimization or refined numerical schemes for accelerated convergence.
    \item The softmax smoothing for preventing empty clusters, while effective \cite{pan2026improved}, introduces a small bias in the mixing proportion estimates.
\end{enumerate}

\subsection{Conclusions}

We have presented a semi-analytically accelerated functional clustering method that combines allometric power-law mean parameterization with SAD(1) covariance modeling. The key methodological contributions are: (1) exploiting the tridiagonal inverse and closed-form determinant of SAD(1) to reduce EM likelihood evaluation from $O(m^3)$ to $O(m)$; (2) deriving closed-form and iteratively closed-form update formulas for $a_{k,i}$ and $\gamma_{k,i}^2$ in the M-step, eliminating general-purpose numerical optimization; and (3) a flexible multivariate extension via block-diagonal covariance. The method is implemented in Python using PyTorch with automatic large-scale initialization, and provides a complete BIC-based model selection workflow.

Future work will explore analytical update formulas for the exponent $b_{k,i}$ (currently frozen in the M-step), integration with network reconstruction frameworks such as idopNetwork \cite{dong2023idopnetwork} for end-to-end clustering-to-network inference, and extensions to higher-order SAD models for long-range dependence structures.

% ============================================================
\section*{Acknowledgments}

The authors thank the contributors of the open-source computational biology community for making their tools and data publicly available.

% ============================================================
\section*{Data Availability}

The \textit{Arabidopsis thaliana} transcriptome data used in this study are publicly available from the Gene Expression Omnibus (GEO). The Python implementation of FunClu-Semi is available at \url{https://github.com/example/funclu-semi}.

\section*{Author Contributions}

{\em To be completed before submission.}

\section*{Competing Interests}

The authors declare no competing interests.

\section*{Funding}

{\em To be completed before submission.}

% ============================================================
% References
% ============================================================
\begin{thebibliography}{99}

\bibitem{xu2005survey}
Xu R, Wunsch D. Survey of clustering algorithms. \textit{IEEE Transactions on Neural Networks}. 2005;16(3):645--678.

\bibitem{banfield1993model}
Banfield JD, Raftery AE. Model-based Gaussian and non-Gaussian clustering. \textit{Biometrics}. 1993;49(3):803--821.

\bibitem{kim2008computational}
Kim BR, Zhang L, Berg A, Fan J, Wu R. A computational approach to the functional clustering of periodic gene-expression profiles. \textit{Genetics}. 2008;180(2):821--834.

\bibitem{wang2012how}
Wang Y, Xu M, Wang Z, Tao M, Zhu J, Wang L, et al. How to cluster gene expression dynamics in response to environmental signals. \textit{Briefings in Bioinformatics}. 2012;13(2):162--174.

\bibitem{wang2014computational}
Wang Y, Wang N, Hao H, Guo Y, Zhen Y, Shi J, et al. A computational algorithm for functional clustering of proteome dynamics during development. \textit{Current Proteomics}. 2014;11(1):1--9.

\bibitem{pan2026improved}
Pan W, Che J, Wu S. An improved method of Wu's functional clustering. \textit{Statistics Innovation}. 2026;3:e005.

\bibitem{west1997general}
West GB, Brown JH, Enquist BJ. A general model for the origin of allometric scaling laws in biology. \textit{Science}. 1997;276(5309):122--126.

\bibitem{huxley1936terminology}
Huxley JS, Teissier G. Terminology of relative growth. \textit{Nature}. 1936;137:780--781.

\bibitem{kingma2014adam}
Kingma DP, Ba J. Adam: a method for stochastic optimization. \textit{arXiv preprint arXiv:1412.6980}. 2014.

\bibitem{nunez1997longitudinal}
N\'{u}\~{n}ez-Ant\'{o}n V. Longitudinal data analysis: non-stationary error structures and antedependent models. \textit{Applied Stochastic Models and Data Analysis}. 1997;13:279--287.

\bibitem{zhao2005non}
Zhao W, Chen YQ, Casella G, Cheverud JM, Wu R. A non-stationary model for functional mapping of complex traits. \textit{Bioinformatics}. 2005;21(10):2469--2477.

\bibitem{zhao2005structured}
Zhao W, Hou W, Littell RC, Wu R. Structured antedependence models for functional mapping of multiple longitudinal traits. \textit{Statistical Applications in Genetics and Molecular Biology}. 2005;4(1):1136.

\bibitem{zimmerman2001parametric}
Zimmerman DL, N\'{u}\~{n}ez-Ant\'{o}n V. Parametric modelling of growth curve data: an overview. \textit{Test}. 2001;10:1--73.

\bibitem{ma2002functional}
Ma CX, Casella G, Wu R. Functional mapping of quantitative trait loci underlying the character process: a theoretical framework. \textit{Genetics}. 2002;161(4):1751--1762.

\bibitem{schwarz1978estimating}
Schwarz G. Estimating the dimension of a model. \textit{The Annals of Statistics}. 1978;6(2):461--464.

\bibitem{weiner2004allocation}
Weiner J. Allocation, plasticity and allometry in plants. \textit{Perspectives in Plant Ecology, Evolution and Systematics}. 2004;6(4):207--215.

\bibitem{sculley2010web}
Sculley D. Web-scale k-means clustering. \textit{Proceedings of the 19th International Conference on World Wide Web}. 2010;1177--1178.

\bibitem{dong2023idopnetwork}
Dong A, Wu S, Che J, Wang Y, Wu R. idopNetwork: a network tool to dissect spatial community ecology. \textit{Methods in Ecology and Evolution}. 2023;14:2272--2283.

\bibitem{james2003clustering}
James GM, Sugar CA. Clustering for sparsely sampled functional data. \textit{Journal of the American Statistical Association}. 2003;98(462):397--408.

\bibitem{kim2010wavelet}
Kim BR, McMurry T, Zhao W, Berg A, Wu R. Wavelet-based functional clustering for patterns of high-dimensional dynamic gene expression. \textit{Journal of Computational Biology}. 2010;17:1067--1080.

\bibitem{cui2008nonparametric}
Cui Y, Zhu J, Wu R. Functional mapping for genetic control of programmed cell death. \textit{Physiological Genomics}. 2008;25:458--469.

\bibitem{li2010functional}
Li N, McMurry T, Berg A, Wang Z, Berceli SA, Wu R. Functional clustering of periodic transcriptional profiles through ARMA($p$,$q$). \textit{PLoS ONE}. 2010;5(4):e9894.

\end{thebibliography}

\end{document}
